\documentclass[11pt, a4paper]{article}
\usepackage[utf8]{inputenc}
\usepackage[UTF8]{ctex} 
\usepackage{amsmath, amssymb, amsfonts}
\usepackage{booktabs} 
\usepackage{graphicx}
\usepackage{geometry}
\usepackage{siunitx}
\usepackage{caption}
\usepackage{subcaption}
\usepackage{float}

\geometry{left=2cm, right=2cm, top=2.5cm, bottom=2.5cm}

\title{\textbf{\part{RLC 谐振电路的稳态特性研究}}}
\author{姓名：党佳一 \quad 学号：20251110109 \\ 北京师范大学物理学系}
\date{实验日期：2026.2.12}

\begin{document}
	
	\maketitle
	
	\section{实验目的}
	\begin{enumerate}
		\item 掌握 RLC 串联谐振电路的稳态特性及谐振频率 $f_0$、品质因数 $Q$ 的测量方法。
		\item 探究电感线圈等效电阻 $R_L$ 随频率变化的非理想物理效应。
		\item 学习使用对数坐标绘制幅频特性与相频特性曲线，验证二阶电路模型。
	\end{enumerate}
	
	\section{实验数据记录与处理}
	
	\subsection{电感线圈等效参数测量}
	使用 LCR 仪在不同频率下对电感 $L$ 及其内阻 $R_L$ 进行测量，结果如表 \ref{tab:LCR} 所示。
	
	\begin{table}[H]
		\centering
		\caption{电感线圈随频率变化的等效参数}
		\label{tab:LCR}
		\begin{tabular}{c S[table-format=5.0] S[table-format=1.3] S[table-format=1.3]}
			\toprule
			序号 & {频率 $f$ (\unit{Hz})} & {电感 $L$ (\unit{mH})} & {等效电阻 $R_L$ (\unit{\Omega})} \\
			\midrule
			1 & 100 & 9.47 & 0.219 \\
			2 & 200 & 9.46 & 0.234 \\
			3 & 500 & 9.46 & 0.225 \\
			4 & 1000 & 9.46 & 0.352 \\
			5 & 2000 & 9.45 & 0.484 \\
			6 & 5000 & 9.446 & 0.924 \\
			7 & 10000 & 9.482 & 1.940 \\
			\bottomrule
		\end{tabular}
	\end{table}
	
	\noindent \textbf{分析：} 电感值在 $100\sim 10000\,\text{Hz}$ 范围内表现非常稳定，基本维持在 $9.45\sim 9.48\,\text{mH}$，说明在该频率范围内，线圈的电感不随频率变化，主要由本身物理性质决定。等效电阻 $R_L$ 随频率 $f$ 显著增长，其主要物理原因包括趋肤效应、邻近效应和磁芯损耗。
	
	\subsection{谐振频率与品质因数的对比测量}
	通过改变 $R, L, C$ 参数，测量不同条件下的谐振特性，对比理论值与实验值。
	
	\begin{table}[H]
		\centering
		\caption{不同电路参数下的谐振点对比测量}
		\begin{tabular}{l c c c}
			\toprule
			参数组 & 理论值 ($f_0 / Q$) & 实验值 ($f_0 / Q$) & 相对误差 \\
			\midrule
			$L=9.481\text{mH}, C=3.3 \times 10^5\text{pF}, R=50\Omega$ & $2.845\text{k} / 3.263$ & $2.84\text{k} / 3.331$ & $0.17\% / 2.0\%$ \\
			$L=9.481\text{mH}, C=1 \times 10^5\text{pF}, R=50\Omega$ & $5.168\text{k} / 5.927$ & $5.2\text{k} / 6.076$ & $0.61\% / 2.5\%$ \\
			$L=9.481\text{mH}, C=1\mu\text{F}, R=50\Omega$ & $1.634\text{k} / 1.874$ & $1.62\text{k} / 1.917$ & $0.76\% / 2.2\%$ \\
			$L=9.481\text{mH}, C=3.3 \times 10^5\text{pF}, R=20\Omega$ & $2.845\text{k} / 7.724$ & $2.838\text{k} / 8.075$ & $0.24\% / 4.5\%$ \\
			\bottomrule
		\end{tabular}
	\end{table}
		\noindent \textbf{分析：} 从数据来看，四组实验谐振频率相对误差均在 $1\%$ 以内，高度验证了 $f_0$ 仅取决于电路的 $L$,$C$参数与回路电阻$R$无关。
		对比第1，4组，可见回路总电阻越小，$Q$ 值越高，符合公式。电阻越小，谐振曲线越尖锐，能耗越低。
		对比1，2，3组实验，$R$不变时，减小$C$，$Q$显著上升，因为减小$C$提高了谐振频率，从而增大感抗。
		$Q$的误差主要来源于$R_L$,因趋肤和邻近效应影响而非定值。
	
	\subsection{频率特性曲线测量数据}
	设置 $R=500\Omega, C=0.1\mu\text{F}, L=9.481\text{mH}$。修正系数 $\kappa = 1.0038$。
	
	\begin{table}[H]
		\centering
		\caption{幅频与相频特性原始数据}
		\begin{tabular}{c S S S S S}
			\toprule
			序号 & {$f$ (kHz)} & {$U$ (V)} & {$U_R$ (V)} & {$\phi$ (度)} & {$\kappa \frac{U_R}{U}$} \\
			\midrule
			1 & 0.1 & 4.02 & 0.125 & -88 & 0.0312 \\
			2 & 0.2 & 4.02 & 0.248 & -85 & 0.0619 \\
			3 & 0.3 & 4.02 & 0.373 & -84 & 0.0932 \\
			4 & 0.4 & 4.07 & 0.498 & -82.0 & 0.1242 \\
			5 & 0.5 & 4.07 & 0.622 & -80.9 & 0.1557 \\
			6 & 0.6 & 3.99 & 0.742 & -79.3 & 0.1887 \\
			7 & 0.7 & 3.98 & 0.859 & -77.5 & 0.2166 \\
			8 & 0.8 & 3.95 & 0.974 & -75.6 & 0.2475 \\
			9 & 0.9 & 3.92 & 1.092 & -74.3 & 0.2796 \\
			10 & 1.0 & 3.90 & 1.209 & -72.5 & 0.3112 \\
			11 & 2.0 & 3.77 & 2.236 & -53.8 & 0.5954 \\
			12 & 3.0 & 3.66 & 3.000 & -35.5 & 0.8228 \\
			13 & 4.0 & 3.59 & 3.438 & -18.72 & 1.0024 \\
			14 & 5.0 & 3.56 & 3.585 & -2.7 & 1.0108 \\
			15 & 6.0 & 3.57 & 3.545 & 10.24 & 0.9968 \\
			16 & 7.0 & 3.59 & 3.402 & 20.73 & 0.9512 \\
			17 & 8.0 & 3.63 & 3.216 & 28.84 & 0.8893 \\
			18 & 25.0 & 3.79 & 1.188 & 57.98 & 0.5451 \\
			19 & 30.0 & 3.89 & 1.039 & 75.1 & 0.2681 \\
			20 & 45.0 & 3.91 & 0.684 & 81.2 & 0.1653 \\
			\bottomrule
		\end{tabular}
	\end{table}
	
	\section{特性曲线拟合结果}
	% 此处用于插入 Python 生成的拟合图
	\begin{figure}[H]
		\centering
		%\includegraphics[width=0.85\textwidth]{C:/RLC frequencycharacteristic.png}
		\caption{RLC 串联电路幅频与相频特性曲线（理论 \&\ 实验）}
		\label{fig:fit}
	\end{figure}
	\section{特性曲线分析与总结}
	
	\subsection{幅频特性分析}
	通过对归一化幅频特性曲线（图\ref{fig:fit}上部）的观察，可以得出以下结论：
	\begin{itemize}
		\item \textbf{谐振峰值与频率：} 实验测量点在 \SI{5.2}{kHz} 处呈现明显的共振峰，与理论值 \SI{5.168}{kHz} 极度接近。归一化增益 $A(\omega)$ 峰值约为 \SI{1.01}{}，验证了谐振时回路总阻抗近似等于总电阻 $R + R_L$ 的物理特性。
		\item \textbf{选择性评价：} 曲线随频率偏离谐振点而迅速衰减。拟合得到的品质因数 $Q \approx 6.08$，表明该系统具有较高的能量存贮效率与频率选择性。
		\item \textbf{对数坐标对称性：} 在对数频率轴上，曲线在 $f_0$ 两侧呈现良好的镜像对称性，这符合二阶线性系统传递函数的数学结构。
	\end{itemize}
	
	\subsection{相频特性分析}
	相频特性曲线（图\ref{fig:fit}下部）清晰地揭示了电路阻抗性质随频率的演变：
	\begin{itemize}
		\item \textbf{相位跨越：} 当 $f < f_0$ 时，相位 $\phi < 0$，电路呈电容性；当 $f > f_0$ 时，$\phi > 0$，电路呈电感性。相位在 $f_0$ 处精确穿过零点，代表此时电压与电流同相。
		\item \textbf{测量灵敏度：} 实验数据点在谐振点附近斜率极高。这种相位的剧烈突变使得利用示波器捕捉“相位差为零”来锁定谐振频率的方法，其精度远高于观察电压幅值极大值的方法。
	\end{itemize}
	
	\subsection{非理想物理效应探讨}
	尽管实验点与理论模型拟合度极高（$R^2 > 0.99$），但细节处仍存在偏离：
	\begin{enumerate}
		\item \textbf{电感等效电阻 $R_L$ 的频散特性：} 实验记录显示 $R_L$ 随频率从 \SI{100}{Hz} 到 \SI{10}{kHz} 增大了约 $9$ 倍。这种由趋肤效应和磁芯损耗引起的 $R_L$ 增加，是导致高频段实验点与使用直流电阻建立的恒定 $Q$ 值理论线产生微小偏差的主要原因。
		\item \textbf{修正系数 $\kappa$ 的有效性：} 通过引入 $\kappa = \frac{R+R_L}{R}$ 对电阻电压进行补偿，有效地将测量量转换为了回路总电阻上的分压，使得实验曲线能够与归一化后的理论包络线高度重合，消除了测量方法带来的系统误差。
	\end{enumerate}
	
	\subsection{结论}
		综上所述，本次实验生成的特性曲线完整地重构了 RLC 串联谐振的物理图景。拟合曲线与实验数据的高度契合，有力地支撑了二阶电路稳态响应理论的正确性，并揭示了真实元件在交流工况下的复杂物理行为。
	\section{一些想法}
	这次试验使我产生了一些关于业余无线电的联想，RLC谐振电路的稳态实验的实质是对信号发生器所发出的特定频率信号的放大，实验二目的是找到谐振电路的谐振点，也就是电路对电信号放大效果最好的频率，品质因数$Q$反应该谐振电路对信号放大效果的特异性，也就是谐振电路放大频率的范围大小，我们且称这个范围为带宽。可以容易得知带宽大小与带宽内频率的放大效果不可兼得，所以$Q$也能够评估谐振电路对带宽内频率的放大效果。实验三的相—频特性曲线也能得出这一点，相位大小正是信号的强弱。当我们把信号发生器换成天线，空间中的携带信息的电磁波驱动着天线内的电荷周期性震荡，馈线将电信号的周期性震荡从天线带入天线调谐器，调谐器中的RLC电路将特定频率的电信号放大，从而便于在下一步的解码中得到原始信息，该分析仅从应用方面考虑，去除了一些精确但又枯燥的数学描述。

	\section{声明}
		本实验报告的$Latex$重构在$Gemini$的参与下完成，但实验数据由本人无$AI$介入下独立完成，实验分析在$AI$协作指导下完成。原始纸质实验报告随信附。
\end{document}